\section{Práce s~ukázkou}

\subsection{Zasazení ukázky do kontextu díla}
Ukázka pochází ze začátku díla. Před ní je krátce představen Švejk a~jeho služebná. Po ní konverzace ještě nějakou dobu pokračuje, než strážník Bretschneider Švejka i~hostinského Plivce zatkne.

\subsection{Časoprostor}
\begin{itemize}
    \item \textbf{Čas:} Děj začíná po zabití Ferdinanda Ferdinanda d'Este což se stalo 28.~června 1914\cite{wiki:Frantisek_Ferdinand} končí přibližně 20.~prosince 1914\cite[174]{osudy_dobreho_vojaka_svejka}
	\item \textbf{Prostor:} Rakousko-Uhersko z~větší části v~Praze
\end{itemize}

\subsection{Téma a~motivy}
\begin{itemize}
    \item \textbf{Téma:} osudy Švejka za světové války
    \item \textbf{Motivy:} válka, úřady, hloupost
\end{itemize}

\subsection{Charakteristika postav}
\begin{itemize}
	\item \textbf{hostinský Palivec} --
	\item \textbf{Švejk} --
\end{itemize}

\subsection{Kompoziční výstavba}
Kompozice je \textbf{chronologická} -- děj na sebe navazuje postupně.

\subsection{Vypravěč a~forma}
Autor je \textbf{vypravěč} -- \textbf{Er forma}, v~průběhu knihy ale Švejk vypráví své historky a~stává se tak vypravěčem v~\textbf{Ich formě}

\subsection{Typy promluv}
Promluvi v knize tvoří dialogy.

\subsection{Veršová výstavba}
Nevyskytuje se, kniha je psána prózou.

\subsection{Jazyková rovina}
Ukázka obsahuje vulgarizmi, spisovnou a mluvenou češtinu a dále se v průběhu knihy pravidelně oběvuje němčina.

\subsection{Umělecké prostředky, tropy a~figury}
\begin{itemize}
    \item \textbf{Kontrast:} „Já se do takových věcí nepletu, s~tím ať mně každej políbí prdel,” odpověděl slušně pan Palivec (...)"; Palivec byl známý sprosťák, (...) Přitom byl ale sečtělý...
    \item \textbf{Gradace:} „Na takovou maličkost se nepamatuju, já jsem se nikdy o~takovou hovadinu nezajímal a~nikdy jsem nebyl na to zvědavej,” (...) „přílišná zvědavost škodí.”
\end{itemize}

\subsection{Vlastní názor na dílo}

\subsection{Filmová a~divadelní zpracování}
\begin{itemize}
    \item \textbf{Muzikál}
    \begin{itemize}
        \item
    \end{itemize}

    \item \textbf{Divadelní adaptace}
    \begin{itemize}
        \item
    \end{itemize}

    \item \textbf{Filmové adaptace}
    \begin{itemize}
        \item
        \item
        \item
    \end{itemize}

\end{itemize}

\newpage

